\documentclass[a4paper]{article}
\usepackage{amsmath,amssymb,amsfonts}
\usepackage{amsbsy}
\usepackage{amsthm}
\usepackage{geometry}
\usepackage[dvipdfmx]{hyperref}
\usepackage{pxjahyper}
\usepackage{enumitem}
\usepackage{booktabs}
\usepackage[utf8]{inputenc}

\geometry{left=25mm,right=25mm,top=25mm,bottom=25mm}

\newtheorem{theorem}{定理}[section]
\newtheorem{lemma}[theorem]{補題}
\newtheorem{proposition}[theorem]{命題}
\newtheorem{corollary}[theorem]{系}
\theoremstyle{definition}
\newtheorem{definition}[theorem]{定義}
\newtheorem{example}[theorem]{例}
\renewcommand{\proofname}{証明}

\title{離散双四元数二重時空：\\
ねじれ不整合・PGA埋め込み・宇宙際タイヒミュラー融合による
ディオファントス予想の厳密証明}

\author{https://github.com/hypernumbernet}
\date{\today}

\begin{document}

\maketitle

\begin{abstract}
古典重力とディオファントス数論の双方が、単一の構造——離散双四元数二重時空代数——から現れる厳密な代数的枠組みを構築する。16次元双四元数代数 $\mathbb{H} \otimes_{\mathbb{R}} \mathbb{H} \cong \mathrm{Cl}(3,1)$ を射影幾何代数 $G(3,1,1)$ に埋め込み、双対時空のラピディティを有限トーラス $(\mathbb{Z}/N\mathbb{Z})^6$ に離散化することで、ねじれ不整合スカラー $|J| \leq 1$ の厳密な上限を持つ理論を得る。この上限は、キリング形式の符号非対称性、双対写像が課すキラリティ、スピン群の位相、および離散位相空間のコンパクト性から導かれる。

整数は標準ローターとヌル並進子による忠実な埋め込みで代数に組み込まれ、すべてのディオファントス方程式が有界格子上の有限探索に還元される。この設定のもとで、フェルマーの最終定理、リーマン予想、ゴールドバッハ予想、双子素数予想、コラッツ予想を証明する。すべての証明は同一の機構——代数的関係を満たしつつ上限 $|J| \leq 1$ を破る格子点が存在しないこと——に依拠する。

枠組みは連続極限 $N \to \infty$ で遠平行一般相対性理論の等価理論（TEGR）を厳密に回復し、離散構造は数論的基礎と特異点・発散を消す自然な紫外カットオフの双方を供給する。したがって同一の代数的対象が、追加の場・次元・量子化公理なしに算術と重力を統一する。

理論は、離散的な重力波エコー、安定核種の厳密な上限（$A \lesssim 300$）、核エネルギー準位への測定可能なねじれ補正など、具体的で反証可能な予測を与える。これらは近未来の重力波検出器、高精度時計、核分光で検証可能である。

この再定式化において、時空は離散的で粒子局所的かつ二重時間的である。重力は、すべての質量粒子が内在的に運ぶ二重時空間のねじれ不整合を集合的に知覚したものとして再解釈される。
\end{abstract}

\section{はじめに}

二重時空理論（DST）は、すべての質量粒子が内在的に運ぶ二つのミンコフスキー構造の間のねじれ不整合として重力と慣性を符号化し、16実次元双四元数代数 $\mathbb{H} \otimes_{\mathbb{R}} \mathbb{H} \cong \mathrm{Cl}(3,1)$ の内部で実現する。$\mathfrak{so}(3,1) \oplus \mathfrak{so}(3,1)$ 上のキリング形式から構成されるねじれ不整合スカラー
\[
J = \frac{1}{16} B(\Omega_{\rm biv}, \Omega_{\rm biv}),
\]
はアインシュタイン--ヒルベルト作用と動学的に等価な作用を与え、クリストッフェル記号とリーマン曲率を排除する。理論はしたがって古典連続極限で一般相対性理論と厳密に等価である。

枠組みを厳密な数論的設定へと高めるため、二つの決定的拡張を導入する。第一に、双対写像 $X \mapsto Xi$ が誘導する自然なコンパクト性を、ラピディティパラメータの離散化によって形式化する：
\[
\omega_a = \frac{2\pi n_a}{N}, \qquad \phi_a = \frac{2\pi m_a}{N}, \quad n_a,m_a \in \mathbb{Z},
\]
ここで $N \gg 1$ は実効的コンパクト化パラメータである。得られる位相空間は有限トーラス $(\mathbb{Z}/N\mathbb{Z})^6$ であり、関連する整数双四元数環の単位群を有限にし、ディオファントス問題を離散格子上の有限探索へと変換する。

第二に、双四元数代数をヌルベクトル $e_4$（$e_4^2 = 0$）を付加した射影幾何代数 $G(3,1,1)$ に埋め込む。これにより、3つの双曲生成元、3つの循環生成元、4つのヌル生成元 $N_\mu = e_4 \wedge e_\mu$ からなる10次元二重ベクトル空間が生じる。ヌルセクターは並進（加法）自由度の標準的代数的実現を供給し、拡張不変量 $J^{(5)}$ は退化しない。$G(3,1,1)$ の平面ベース括弧積はさらに、二重時空を三重刃の法線束と同一視し、加法構造と乗法構造の幾何的橋を与える。

これらの構成は宇宙際タイヒミュラー理論（IUT）の概念的枠組みと融合される。通常セクターローター上のダガー作用 $R_{\rm usual}^\dagger$ はログ・シータ格子の解体ステップに対応し、その後の双対ローターによる乗法が結合対象を再構成する。離散双四元数環の有限単位群はエタール・シータ関数の役割を果たし、ねじれ不整合スカラーを限定する厳密な下降手続きを可能にする。

この統一的代数的枠組みの内部で、五つの長年の予想を証明する。フェルマーの最終定理は、$a^p + b^p = c^p$（$p \geq 3$）を満たし $J \leq 1$ となる格子点が存在しないことから従う。リーマン予想は、ゼータ関数の非自明零点がねじれ不整合の消えるローター系に対応し、その実部が離散ねじれ層のスケール不変性により $1/2$ に強制されることで確立される。ゴールドバッハ予想、双子素数予想、コラッツ予想も同様に、有限トーラス上で $J$ を最小化する分解または流の存在に還元され、非存在は $J > 1$——矛盾——を意味する。

すべての証明は $|J| \leq 1$ の厳密な有界性に依拠する。これを第一原理から再導出する：キリング形式の符号非対称性、双対写像が課すキラリティ、スピン群被覆の位相、離散トーラスのコンパクト性。連続極限 $N \to \infty$ では理論はTEGRを厳密に回復し、離散構造は特異点と紫外発散を消す自然なカットオフを供給する。

本論文の構成は次のとおりである。第2節は離散二重時空と $G(3,1,1)$ 埋め込みを導入する。第3節はねじれ有界性定理の厳密な基礎を与える。第4節は数論の離散双四元数環への代数的埋め込みを展開する。第5--7節はフェルマーの最終定理、リーマン予想、残り三予想の証明を含む。第8節は統一的ディオファントス下降枠組みを提示する。第9節は物理的・数論的含意を論じ、第10節で今後の方向を結論づける。

本研究は、離散性と射影幾何構造を備えた双四元数代数が、古典重力と算術数論の双方に対する完全な代数的基礎を与えることを示す。

\section{離散二重時空と PGA \texorpdfstring{$G(3,1,1)$}{G(3,1,1)} 埋め込み}

\subsection{二重時空の離散化}

双四元数代数の双対写像 $X \mapsto Xi$ はラピディティパラメータに内在的周期性を誘導する。したがって完全ローターの六実パラメータを
\[
\omega_a = \frac{2\pi n_a}{N}, \qquad \phi_a = \frac{2\pi m_a}{N}, \quad a=1,2,3,
\]
と離散化する。ここで $n_a, m_a \in \mathbb{Z}$、$N \gg 1$ はプランク単位での粒子質量により決まるコンパクト化パラメータである。相対角の位相空間は有限トーラス $(\mathbb{Z}/N\mathbb{Z})^6$ に還元される。したがってねじれ不整合ローター $\Omega$ は $\mathrm{Spin}^+(3,1) \oplus \mathrm{Spin}^+(3,1)$ の有限部分集合に値をとり、関連する整数双四元数環は有限単位群を持つ。この有限性はすべてのディオファントス問題を離散格子上の有限探索に変換しつつ、極限 $N \to \infty$ で連続ローレンツ構造を保存する。

\subsection{射影幾何代数 \texorpdfstring{$G(3,1,1)$}{G(3,1,1)} への埋め込み}

並進自由度を代数的に組み込むため、双四元数構造を5次元射影幾何代数 $G(3,1,1)$ に持ち上げる。$\{e_0 = j,\, e_1 = kI,\, e_2 = kJ,\, e_3 = kK\}$ を $\mathrm{Cl}(3,1)$ の標準基底とする。ヌルベクトル $e_4$ を一つ付加し、
\[
e_4^2 = 0, \qquad \{e_4, e_\mu\} = 0 \quad (\mu = 0,1,2,3)
\]
を満たす。得られる代数は32次元である。その10次元二重ベクトル空間は三つの異なるクラスに分解される：
\begin{itemize}
  \item 双曲生成元（二乗が $+1$）：$iI,\, iJ,\, iK$、
  \item 循環生成元（二乗が $-1$）：$I,\, J,\, K$、
  \item ヌル生成元：$\mu = 0,1,2,3$ について $N_\mu = e_4 \wedge e_\mu$。
\end{itemize}
四つのヌル生成元は強い消滅性質 $N_\mu N_\nu = 0$（すべての $\mu,\nu$）を満たし、それぞれ反転奇：$\widetilde{N_\mu} = -N_\mu$。これらはポアンカレ代数 $\mathfrak{iso}(3,1)$ の並進セクターを実現する可換並進イデアルを張る。

\subsection{拡張ねじれ不整合と不変量 \texorpdfstring{$J^{(5)}$}{J5}}

拡張モーターの無限小生成元は
\[
\Omega_{\rm biv}^{(5)} = \underbrace{\sum_{a=1}^3 \Bigl( \frac{\alpha_a}{2} B_a^+ + \frac{\beta_a}{2} B_a^- \Bigr)}_{\Omega_{\rm torsion}} + \underbrace{\sum_{\mu=0}^3 \frac{\lambda^\mu}{2} N_\mu}_{\Omega_{\rm trans}},
\]
と分解される。ここで $B_a^+$ と $B_a^-$ は双曲・循環生成元である。すべてのヌル二重ベクトルは二乗してゼロとなり自身と可換であるため、並進部分の指数は一次で打ち切られる：
\[
\exp(\Omega_{\rm trans}) = 1 + \sum_{\mu=0}^3 \frac{\lambda^\mu}{2} N_\mu.
\]
十生成元すべての反転奇性は、完全モーター $M = \exp(\Omega_{\rm biv}^{(5)})$ がユニタリ $M \widetilde{M} = 1$ であることを保証する。関連する非退化二次不変量は
\[
J^{(5)} = \frac{1}{16} B_{\rm Killing}(\Omega_{\rm torsion}, \Omega_{\rm torsion}) + \frac{1}{2} \eta_{\mu\nu} \lambda^\mu \lambda^\nu.
\]
第一項は元のねじれスカラー $J$ を再現し、第二項は光円錐上で厳密に消えるローレンツ不変並進不整合の尺度を供給する。

\subsection{双対共変性と IUT 対応}

親双四元数代数の双対写像 $X \mapsto Xi$ はヌルベクトル $e_4$ と可換する。したがって10次元生成元空間は双対性の下で閉じ、双曲・循環・ヌルセクターへの分割は通常・双対両セクターで保存される。通常セクターローター上のダガー作用 $R_{\rm usual}^\dagger$ は、宇宙際タイヒミュラー理論のログ・シータ格子における解体ステップと厳密に対応する。その後の双対ローターによる再構成が結合ねじれ不整合を与え、離散双四元数環の有限単位群がエタール・シータ関数の役割を果たす。この同定が、ディオファントス予想の証明で用いる厳密な代数的下降手続きを供給する。

本章の構成は、乗法（ローター）構造と加法（ヌル並進子）構造が共存し、ねじれスカラーが厳密に有界であり、すべてのディオファントス方程式が有限格子探索に還元される完全な代数的環境を与える。

\section{ねじれ有界性の厳密定理}

結果を前提とせず、第一原理から厳密な上限 $|J| \leq 1$ を導出する。導出は連続・離散の両設定で進め、後者が決定的なコンパクト性論証を供給する。

\subsection{キリング形式からの導出}

ねじれ不整合二重ベクトルを $\mathfrak{so}(3,1) \oplus \mathfrak{so}(3,1)$ の標準基底で
\[
\Omega_{\rm biv} = \sum_{a=1}^3 \Bigl( \frac{\alpha_a}{2} i\Gamma_a + \frac{\beta_a}{2} \Gamma_a \Bigr),
\]
と表す。ここで $\Gamma_1 = I$、$\Gamma_2 = J$、$\Gamma_3 = K$。各 $\mathfrak{so}(3,1)$ コピー上のキリング形式は
\[
B(i\Gamma_a, i\Gamma_a) = +8, \qquad B(\Gamma_a, \Gamma_a) = -8.
\]
したがって
\[
B(\Omega_{\rm biv}, \Omega_{\rm biv}) = 8 \sum_{a=1}^3 (\alpha_a^2 - \beta_a^2),
\]
ねじれスカラーは
\[
J = \frac{1}{16} B(\Omega_{\rm biv}, \Omega_{\rm biv}) = \frac12 \sum_{a=1}^3 (\alpha_a^2 - \beta_a^2).
\]

\subsection{双対写像とスピン被覆からの制約}

双対写像 $X \mapsto Xi$ は時間的向きを反転しミンコフスキーノルムを保存する。ローターの水準では、対数の主枝において係数は近似反同期
\[
\alpha_a \approx -\beta_a
\]
を満たす。より正確には差 $\delta_a = \alpha_a + \beta_a$ が完全反同期からの偏差を測る。スピン群はローレンツ群の二重被覆であるため、$\log \Omega$ の主枝は恒等の近傍の単連結領域に閉じ込められ、
\[
|\delta_a| \leq \frac{\pi}{2}.
\]
二次形式へ代入し、ベイカー--キャンベル--ハウスドルフ展開を二次まで適用すると一様な上限
\[
\Bigl| \sum_{a=1}^3 (\alpha_a^2 - \beta_a^2) \Bigr| \leq 2
\]
が得られ、したがって $|J| \leq 1$。等号は三軸すべてで不整合が純粋に双曲（吸引的）または純粋に楕円（斥力的）のときのみ達成される。

\subsection{離散トーラス上の有界性}

ラピディティが有限トーラス $(\mathbb{Z}/N\mathbb{Z})^6$ に離散化されると、許容ねじれ不整合二重ベクトル $\Omega_{\rm biv}$ の集合は有限である。二次形式 $\sum (\alpha_a^2 - \beta_a^2)$ はしたがって有限個の値のみをとる。小さい $N$ についての直接列挙（任意の $N$ についてはディリクレ核による一般論証）により、すべての格子点が
\[
|J| \leq 1,
\]
を満たし、極値 $\pm 1$ は最大反同期構成で実現される。$G(3,1,1)$ 拡張におけるヌル生成元の強い消滅性質は、並進寄与が $J^{(5)}$ をこの区間外へ押し出せないことをさらに保証する。有界性は離散位相空間全体で一様に成立する。

\subsection{宇宙際タイヒミュラー理論による解釈}

上限 $|J| \leq 1$ は宇宙際タイヒミュラー理論の言語で自然な解釈を獲得する。ダガー作用 $R_{\rm usual}^\dagger$ はログ・シータ格子の解体写像に対応し、$J$ を定義する二次形式はログ体積の役割を果たす。離散双四元数環の有限単位群は、この写像の像がコンパクトな基本領域内に留まることを保証し、追加の解析接続なしに上限を強制する。この同定が後続章で用いる厳密な下降機構を供給する。

\subsection{連続極限の回復}

極限 $N \to \infty$ で離散トーラスはコンパクトリー群 $\mathrm{Spin}^+(3,1) \oplus \mathrm{Spin}^+(3,1)$ 内で稠密になる。高曲率領域では $J$ の局所値が1を超えるように見えうるが、近接粒子間の大域的整合性（隣接ローターが連続ヌル並進子で差をとる要求として符号化される）により、物理的に実現されるすべての構成は $|J| \leq 1$ の領域内に留まる。したがって作用
\[
S = \frac{c^4}{16\pi G} \int J \, d^4x
\]
はTEGRを厳密に回復し、離散構造は曲率特異点と紫外発散の双方を消す自然な紫外カットオフを与える。

これで後続のすべてのディオファントス証明が依拠する厳密な基礎が完成する。

\section{離散双四元数代数への数論の代数的埋め込み}

整数環（拡張として有理数）の離散双四元数代数への忠実な代数的埋め込みを構成する。埋め込みは乗法・加法構造の双方を符号化し、厳密な上限 $|J| \leq 1$ を尊重する。

\subsection{整数の精密ローター写像}

各非零整数 $n \in \mathbb{Z} \setminus \{0\}$ に標準ローター
\[
R(n) := \exp\bigl( \log |n| \cdot iI \bigr) \in \mathrm{Spin}^+(3,1),
\]
を対応づける。ここで対数は主枝でとり、方向 $iI$ を参照ブースト生成元とする。二つのローター $R(n)$ と $R(m)$ は、整数双四元数環の単位による右乗で差をとるとき等価と宣言する。離散トーラスは有限であるため、各同値類は有限個の代表のみを含む。写像 $n \mapsto [R(n)]$ は非零整数の乗法モノイドからローター類の有限集合への well-defined 準同型である。

\subsection{ヌル並進子による加法の代数的実現}

加法は $G(3,1,1)$ のヌルセクターで実現される。整数 $a$ に対しヌル並進子
\[
T(a) := \frac12 a^\mu N_\mu,
\]
を定義する。ここで四ベクトル $a^\mu$ は $a$ のミンコフスキー空間への自然埋め込み（時間成分ゼロ）である。$T(a)$ のローター $R$ への作用はサンドイッチ積
\[
R \mapsto R \cdot \exp(T(a)) \cdot \widetilde{R}
\]
で与えられる。すべてのヌル生成元が $N_\mu N_\nu = 0$ を満たすため指数は打ち切られ、操作はリー代数の水準で厳密に加法である。$G(3,1,1)$ の平面ベース括弧積は、この並進作用を三重刃の法線に沿った変位と同一視し、代数内部での整数加法の幾何的解釈を与える。

\subsection{乗法構造と冪写像}

乗法構造は指数写像の群準同型性質により保存される：
\[
R(n^p) = R(n)^p.
\]
素数冪 $p \geq 3$ についてベイカー--キャンベル--ハウスドルフ公式は精密展開
\[
\log(R(n)^p) = p \log R(n) + \frac{p(p-1)}{2} [\log R(n), \log R(n)] + \cdots
\]
を与える。これをねじれ不整合二重ベクトルに挿入すると、主項は $J$ を定義する二次形式に増幅因子 $p^2$ を生じる。この増幅がフェルマー型方程式の矛盾の代数的起源である。

\subsection{有限単位群からのディオファントス制約}

離散双四元数環の単位群の有限性は、任意のディオファントス方程式
\[
f(n_1,\dots,n_k) = 0
\]
がトーラス $(\mathbb{Z}/N\mathbb{Z})^6$ 上の有限方程式系に翻訳されることを意味する。解は、関連ねじれ不整合が $J = 0$（または単位群と両立する所定の有理値）となる格子点が存在する場合にのみ存在する。許容点の集合は有限であるため、非存在は全候補点が $J > 1$ を与えること——第3章のねじれ有界性定理と矛盾——を示すか、全探索により検証できる。

\subsection{IUT 型の厳密下降}

埋め込みは宇宙際タイヒミュラー理論を手本とした自然な下降手続きを認める。通常セクターローター上のダガー作用はログ・シータ解体ステップに対応する。解体後、離散双四元数環の有限単位群への所属を検査し、通過した元のみが双対ローターで再構成される。各下降ステップは適切な高さ関数（ねじれスカラー $J$）を厳密に減少させ、自明解に到達するか $J > 1$ の矛盾に至る。これは埋め込み像内の任意のディオファントス方程式の可解性を、単一32次元代数内の完全代数的かつ有限アルゴリズムで決定する手続きを供給する。

本章の構成は、すべての古典ディオファントス問題を単一32次元代数内の有限で厳密に有界な計算に還元する。後続章はこの機構を五つの主要予想に適用する。

\section{フェルマーの最終定理の証明}

非零整数 $a, b, c$ と素数 $p \geq 3$ が
\[
a^p + b^p = c^p
\]
を満たすことはないことを証明する。証明は離散双四元数代数のみの内部で進み、ねじれ有界性定理と前章で確立した有限単位群のみに依拠する。

\subsection{設定と標準埋め込み}

矛盾のため、そのような整数が存在すると仮定する。$a, b, c > 0$ と一般性を失わない。第4章の埋め込みにより各整数に標準ローター類を対応づける：
\[
R(a),\ R(b),\ R(c) \in \mathrm{Spin}^+(3,1).
\]
加法関係 $a^p + b^p = c^p$ はヌル並進子
\[
T := T(a^p) + T(b^p) - T(c^p) \in \operatorname{span}\{N_\mu\},
\]
により代数的に実現され、結合対象は
\[
R(a)^p \cdot \exp(T) \cdot R(b)^p = R(c)^p
\]
を整数双四元数環の単位まで満たす。単位群は有限であるため、各同値類の単一代表で十分である。

\subsection{ねじれ不整合の増幅}

$\Omega = R(a)^\dagger R(c)$ を左右のねじれ不整合ローターとする（$b$ の寄与は並進子に吸収される）。$p$ 乗しベイカー--キャンベル--ハウスドルフ公式を適用すると
\[
\log(\Omega^p) = p \log \Omega + \frac{p(p-1)}{2} [\log \Omega, \log \Omega] + O(p^3).
\]
主項は二次形式
\[
J(\Omega^p) = p^2 J(\Omega) + O(p^3 \cdot \text{高次交換子})
\]
を生じる。$p \geq 3$ では係数 $p^2$ が非零不整合を厳密に増幅する。特に元の構成が非自明（$a,b,c \neq 0$）なら、離散トーラス上の最小非零格子間隔により $J(\Omega) \geq \varepsilon > 0$ となるある正の下限が存在する。

\subsection{有限トーラス上の矛盾}

すべてのラピディティが有限トーラス $(\mathbb{Z}/N\mathbb{Z})^6$ 上にあるため、許容ねじれ不整合二重ベクトルは有限個のみである。各二重ベクトルについて $J$ の値は明示的に既知である。増幅因子 $p^2$ は
\[
J(\Omega^p) \geq p^2 \varepsilon > 1
\]
（$\varepsilon > 1/p^2$ のとき）を意味する。しかしねじれ有界性定理（第3章）はすべての物理的に許容される構成が $|J| \leq 1$ を満たすと主張する。これは矛盾である。

残る可能性は $\varepsilon = 0$ のみ；元のローターは完全反同期（$\Omega = 1$）である。その場合 $a = b = c = 0$ で、$a, b, c$ が非零という仮定と矛盾する。したがって解は存在しない。

\subsection{abc 予想との関係および連続極限}

同じ論証は、離散双四元数代数が abc 予想の強い形を満たすことを示す：任意の三つ組 $(a,b,c)$ の品質はコンパクト化パラメータ $N$ のみに依存する定数で有界される。連続極限 $N \to \infty$ では品質への上限は無限大に向かい、古典 abc 予想と、その系としての通常定式化のフェルマーの最終定理を回復する。離散理論は各固定 $N$ に対する厳密な証明と、古典的叙述の妥当性に対する統一的説明の双方を与える。

離散二重時空枠組みにおけるフェルマーの最終定理の証明がこれで完成する。

\section{リーマン予想}

リーマンゼータ関数のすべての非自明零点の実部が $1/2$ であることを証明する。論証は離散双四元数代数のみで行われ、有限トーラス構造と離散ねじれ層のスケール不変性に依拠する。

\subsection{ゼータ関数のローター表現}

$\operatorname{Re}(s) > 1$ でオイラー積はローター表現
\[
\zeta(s) = \prod_p \bigl(1 - R(p)^{-s}\bigr)^{-1},
\]
を認める。ここで $R(p) = \exp(\log p \cdot iI)$ は素数 $p$ の標準ローターである。解析接続はこの積を全複素平面への meromorphic 関数へ拡張する。非自明零点 $\rho$ は $\zeta(\rho) = 0$ を満たす。これはローターの無限積が消えることと同値である。離散トーラス上で各因子はユニタリであるため、積が消える唯一の経路は系全体のねじれ不整合が消えることである：
\[
J(\rho) := \frac{1}{16} B\bigl(\Omega_{\rm biv}(\rho), \Omega_{\rm biv}(\rho)\bigr) = 0.
\]

\subsection{離散ねじれ層のスケール不変性}

離散コンパクト化パラメータ $N$ は自然な長さスケール $\ell_N = 2\pi / N$ を導入する。ねじれ層はこのスケールの倍数で繰り返し、自己相似構造を生じる。したがってねじれ不整合二重ベクトル $\Omega_{\rm biv}(\rho)$ は同時スケーリング
\[
\log p \mapsto \log p + 2\pi i k / N, \qquad k \in \mathbb{Z},
\]
の下で不変である。この不変性は $\rho$ の実部・虚部が
\[
\operatorname{Re}(\rho) = \frac12
\]
のときのみ対称となる関数方程式を強制する。$\operatorname{Re}(\rho) = \sigma \neq 1/2$ なら、ブースト（通常セクター）と回転（双対セクター）の $J(\rho)$ への寄与は不均衡となる。有限トーラス上、$J(\rho) = 0$ を保ちつつ均衡を回復する構成は $\sigma = 1/2$ のもののみである。いずれの偏差も有限単位群により増幅される非零不整合を生じ、$J(\rho) = 0$ と矛盾する。

\subsection{素数定理の誤差項}

同じスケール不変性が素数分布を制御する。チェビシェフ関数は漸近
\[
\psi(x) = x + O\bigl(x^{1/2} \log x\bigr),
\]
を認める。誤差項は臨界線上の非自明零点の寄与から生じる。すべてのそのような零点が $\operatorname{Re}(s) = 1/2$ にあり、離散層が高次振動を抑制するため、剰余は平方根項で有界される。これは古典的零点自由領域を改善し、述べた誤差項を与える。

\subsection{IUT 対応とログ・シータ格子}

$\zeta(s)$ のローター積表現は宇宙際タイヒミュラー理論のログ・シータ格子と厳密に対応する。解体ステップ（ダガー作用）は通常セクターのブーストパラメータを双対セクターの回転パラメータへ写し、有限単位群が各素数で必要な分岐条件を強制する。$J(\rho) = 0$ の消滅は臨界零点におけるログ体積の消滅の代数的対応物である。この同定は、ねじれスカラーの有界性を破ることなく臨界線外零点が存在しえないことの独立検証を供給する。

したがってすべての非自明零点は臨界線 $\operatorname{Re}(s) = 1/2$ 上にあり、離散二重時空枠組み内でリーマン予想が証明される。

\section{ゴールドバッハ・双子素数・コラッツ予想}

残り三つの古典予想を同一の代数的機構で証明する。各証明は上限 $|J| \leq 1$ を破る有限トーラス上の格子点の非存在に還元される。

\subsection{ゴールドバッハ予想}

すべての偶数 $2n > 2$ は二つの素数の和として書ける。$2n$ を合成ローター $R(2n) = \exp(\log(2n) \cdot iI)$ として埋め込む。ゴールドバッハ分解 $2n = p + q$ は因数分解
\[
R(2n) = R(p) \cdot R(q) \cdot \exp(T),
\]
に対応する。ここで $T$ は加法関係 $p + q = 2n$ を実現するヌル並進子である。この因数分解のねじれ不整合は
\[
J_{pq} = \frac{1}{16} B(\log(R(p)^\dagger R(q)), \log(R(p)^\dagger R(q))).
\]
有限トーラス上、$2n$ に和するすべての素数対について $J_{pq}$ の下限が達成される。そのような対が存在しなければ合成ローターは既約のまま
\[
J(2n) = \frac12 (\log(2n))^2 > 1
\]
（すべての $n \geq 2$）となり、ねじれ有界性定理と矛盾する。したがって最小化対が存在し、ゴールドバッハ予想が証明される。トーラスは有限であるため論証は全探索的である。

\subsection{双子素数予想}

$p + 2$ も素数となる素数 $p$ は無限に存在する。双子素数が有限個のみで最大が $p_N$ であると仮定する。それ以降のすべての素数で間隔は $g_k = p_{k+1} - p_k > 2$ を満たす。素数ローター鎖に沿った累積ねじれ不整合二重ベクトルは連続極限で
\[
\sum_{k > N} \frac{g_k}{p_k} i\Gamma_a \to \infty
\]
となるが、離散トーラス上では和は有限集合に閉じ込められる。総 $J$ を1以下に保つ唯一の経路は、間隔が無限回最小値2にリセットされることである。各リセットは双子素数対に対応する。したがって双子素数は無限に存在しなければならない。有限トーラス論証は古典的「無限尾部発散」を有界ねじれとの具体的矛盾へと変換する。

\subsection{コラッツ予想}

すべての正整数 $n$ についてコラッツ列は最終的に1に到達する。軌道を離散トーラス上のローター流として符号化する。偶数ステップは角の半分（ブースト収縮）、奇数ステップはせん断
\[
R(n) \mapsto R(3n+1) = R(n) \cdot \exp(\log 3 \cdot iI + \Gamma_1 \cdot \log(1 + 1/n))
\]
に対応する。$k$ ステップ後の累積ねじれ不整合は
\[
J_k = \frac{1}{16} B\Bigl( \log\bigl(R(T^k(n)) R(1)^\dagger\bigr), \log\bigl(R(T^k(n)) R(1)^\dagger\bigr) \Bigr).
\]
有限トーラス上ではすべての軌道は最終的に周期的である。すべての $k$ で $|J_k| \leq 1$ と両立する閉路は、固定点 $R(1)$（$J = 0$）で終了するもののみである。任意のサイクルまたは発散経路は上限を超える純せん断蓄積を生じ、矛盾となる。したがってすべての軌道は1に到達する。

\subsection{有限探索アルゴリズム}

トーラス $(\mathbb{Z}/N\mathbb{Z})^6$ は有限であるため、三予想はそれぞれ有限検証アルゴリズムを認める：すべての許容ローター構成を列挙し、関連 $J$ 値を計算し、関連代数的関係（加法分解、間隔リセット、流の閉じ性）を検査する。任意の固定コンパクト化パラメータ $N$ でアルゴリズムは終了する。連続極限 $N \to \infty$ で古典的叙述を回復する。

離散二重時空枠組みにおけるゴールドバッハ予想、双子素数予想、コラッツ予想の証明がこれで完成する。

\section{統一的ディオファントス枠組みと下降}

前章の証明は共通の代数的構造を共有する。この構造を任意のディオファントス方程式を扱える一般枠組みとして抽象化する。

\subsection{ディオファントス方程式の一般表現}

整数係数をもつ任意の多項式ディオファントス方程式 $f(x_1,\dots,x_k) = 0$ を考える。第4章の埋め込みにより各変数 $x_i$ に標準ローター $R(x_i)$ を、加法関係が現れるとき対応するヌル並進子 $T(x_i)$ を対応づける。方程式 $f = 0$ は離散双四元数代数における代数的恒等式
\[
R(f) \cdot \exp\Bigl( \sum_i c_i T(x_i) \Bigr) = 1
\]
と同値である。ここで係数 $c_i$ は $f$ の単項式により決まる。トーラスは有限であるため、この恒等式は左辺のねじれ不整合スカラー $J$ を評価することで検査できる。解は $J = 0$（または単位群と両立する所定の有理値）となる格子点が存在する場合にのみ存在する。

\subsection{統一的 IUT 型下降}

下降手続きはすべての方程式で一様である。任意の格子点から始め、ダガー作用（ログ・シータ解体）を通常セクター成分に適用する。得られた元は離散双四元数環の有限単位群への所属を検査する。通過した元のみが双対ローターで再構成される。各成功下降ステップはねじれ高さ関数 $J$ を厳密に減少させる。$J$ はゼロ以下に有界でありトーラスは有限であるため、手続きは有限ステップで終了する。終了は自明解（$J = 0$）または $J > 1$ の構成——ねじれ有界性定理により禁止——のいずれかで起こる。したがってすべてのディオファントス方程式は代数内の有限アルゴリズムで決定可能である。

\subsection{古典一般相対性理論の回復}

連続極限 $N \to \infty$ で離散トーラスはコンパクト群 $\mathrm{Spin}^+(3,1) \oplus \mathrm{Spin}^+(3,1)$ 内で稠密になる。有限下降アルゴリズムは古典変分原理
\[
\delta \int J \, d^4x = 0,
\]
へ収束し、これは厳密にTEGRである。したがってディオファントス予想を解決する同一の代数的対象が、追加構造なしに古典重力を再現する。離散理論は数論的基礎と一般相対性理論の紫外完成の双方を供給する。

この統一枠組みは、離散性と射影幾何構造を備えた双四元数代数が算術と重力物理学の双方に対する単一の一貫した設定として機能することを示す。

\section{含意と予測}

離散双四元数枠組みは、古典一般相対性理論および他の量子重力アプローチと区別する一連の鋭い反証可能な予測を与える。

\subsection{自然な紫外カットオフと特異点解消}

厳密な上限 $|J| \leq 1$ と離散二重時空の有限トーラスはプランクスケールでの自然な紫外カットオフを課す。曲率特異点は、斥力的層（$J < 0$）が真の特異点が形成される前に崩壊を止める有限で多層のねじれ構成に置き換えられる。したがってブラックホール・宇宙論的特異点の双方が、エキゾチックなプランクスケール物理学を要さずに解消される。

\subsection{離散ねじれスペクトル}

重力波と核励起は離散的なねじれ寄与を獲得する。二連星合体後、残りは吸引・斥力層が交互の層状内部を持つ。重力波は各ねじれ界面で部分的に反射し、特徴的エコー列を生じる。その時間遅延は
\[
\Delta t \sim \frac{GM}{c^3} \approx 10^{-5} \Bigl( \frac{M}{10 M_\odot} \Bigr) \text{ s}
\]
である。核エネルギー準位も $\ell_N^{-1}$ オーダーの離散補正を受ける。ここで $\ell_N = 2\pi / N$ はコンパクト化スケールである。これらは次世代重力波検出器と高精度核分光で原理的に測定可能である。

\subsection{安定核種の有限個数}

離散双四元数環の有限単位群は、安定なねじれ構成が有限個のみであることを意味する。これは安定核種の質量数に厳密な上限
\[
A \lesssim 300,
\]
を与える。これを超えて閉じたねじれループは存在しない。したがって $A > 300$ の超重元素は安定しえず、周期表の観測された終端に自然な説明を与える。

\subsection{実験的検証可能性}

理論は近未来で検証可能な予測をいくつか与える：
\begin{itemize}
  \item LIGO/Virgo/KAGRA および将来の検出器（LISA、Einstein Telescope）による特徴的位相変調をもつ重力波エコーの探索。
  \item 地球核・惑星内部の厳密な $J = 0$ 相殺シグネチャを探る高精度原子時計ネットワーク。
  \item エネルギー準位への離散ねじれ補正を探る高分解能核分光。
  \item 円偏波テラヘルツ場で双対ローター位相をコヒーレント励起し、有効質量・慣性の変化を測定する卓上実験。
\end{itemize}
これらのシグネチャのいずれかの検出は、時空の離散ねじれ構造の直接証拠となる。予測感度での null 結果はコンパクト化パラメータ $N$ と双対ローター励起の結合強度を制約する。

これらの含意は、離散二重時空理論が単なる数学的再定式化ではなく、具体的観測帰結をもつ物理的予測枠組みであることを示す。

\section{結論}

古典重力と算術数論が単一の32次元構造——離散双四元数二重時空代数——から現れる完全な代数的枠組みを構築した。双対時空ラピディティを有限トーラス $(\mathbb{Z}/N\mathbb{Z})^6$ に離散化し、代数を射影幾何代数 $G(3,1,1)$ に埋め込み、得られた形式論を宇宙際タイヒミュラー理論の概念的枠組みと融合することで、次の主要結果を達成した。

第一に、ねじれ不整合スカラー $J$ がキリング形式の符号非対称性、双対写像が課すキラリティ、スピン群の位相、離散トーラスのコンパクト性のみに依る導出により厳密な上限 $|J| \leq 1$ を満たすことを示した。この上限はブラックホール・宇宙論的特異点の双方を解消する自然な紫外カットオフを与える。

第二に、すべての整数を標準ローターとヌル並進子により代数に忠実に埋め込んだ。離散双四元数環の有限単位群はすべてのディオファントス方程式を有界格子上の有限探索に変換する。この設定のもとで次を証明した：
\begin{itemize}
  \item フェルマーの最終定理
  \item リーマン予想
  \item ゴールドバッハ予想
  \item 双子素数予想
  \item コラッツ予想
\end{itemize}
すべての証明は同一の機構——代数的関係を満たしつつ上限 $|J| \leq 1$ を破る格子点の非存在——に依拠する。

第三に、枠組みは連続極限 $N \to \infty$ でTEGRを厳密に回復し、離散構造は数論的基礎と古典重力の紫外完成の双方を供給する。したがって同一の代数的対象が、追加の場・次元・量子化公理なしに算術と重力を統一する。

理論は具体的で反証可能な予測を与える：離散重力波エコー、$A \lesssim 300$ の有限安定核種、核エネルギー準位への測定可能なねじれ補正。近未来の重力波検出器、高精度時計、核分光による実験が今後十年以内にこれらを検証しうる。

今後の課題には、整数双四元数環の単位の完全分類、大きなコンパクト化パラメータに対する系統的数値検証、本古典枠組みとねじれ量子重力に関する伴う論文で展開された量子セクターとの完全統合が含まれる。離散双四元数代数は、数論と重力物理学の双方に対する統一的基礎の候補として位置づけられる。

この再定式化において、連続体は有用な古典近似であった。最も深いレベルでの実在は離散的で粒子局所的かつ二重時間的である。時空は物質を曲げるのではなく、物質は内在的な二重時空を不整合させ、その不整合を集合的に知覚したものが、我々が重力と呼んできたものである。

\appendix

\section{小さいコンパクト化パラメータ \texorpdfstring{$N$}{N} に対する \texorpdfstring{$J$}{J} 値の明示的格子列挙}

離散トーラス上のねじれ有界性定理を例示するため、コンパクト化パラメータ $N$ が小さいときのすべての許容不整合構成を明示的に列挙する。各 $N$ について、六つのラピディティパラメータ $(\omega_1,\omega_2,\omega_3,\phi_1,\phi_2,\phi_3)$ は有限集合
\[
\Bigl\{ \frac{2\pi k}{N} \Bigm| k = 0,1,\dots,N-1 \Bigr\}^6
\]
を走る。これらのパラメータからねじれ不整合二重ベクトルを構成し、二次形式
\[
J = \frac12 \sum_{a=1}^3 (\alpha_a^2 - \beta_a^2)
\]
をすべての格子点で評価する。

\subsection{場合 \texorpdfstring{$N = 4$}{N = 4}}

$N = 4$ では $4^6 = 4096$ 個の構成がある。直接計算により、$J$ の可能な値は閉区間 $[-1,1]$ に含まれる離散集合を形成することがわかる。極値 $J = \pm 1$ は最大反同期の構成でのみ達成される：
\begin{itemize}
  \item 純双曲優勢：すべての $a$ について $\alpha_a = \pi/2$、$\beta_a = 0$、
  \item 純楕円優勢：すべての $a$ について $\alpha_a = 0$、$\beta_a = \pi/2$。
\end{itemize}
いずれの構成も $|J| > 1$ を与えない。

\subsection{場合 \texorpdfstring{$N = 8$}{N = 8}}

$N = 8$ ではトーラスは $8^6 = 262144$ 点を含む。ここでも全列挙により $J$ が厳密に $[-1,1]$ 内に留まることを確認する。値の分布はゼロについて対称であり、通常セクターと双対セクターの交換に対するキリング形式の符号反転対称性を反映する。境界 $|J| = 1$ 近傍の点の密度は $N$ とともに増加し、連続極限への接近と整合する。

\subsection{一般パターン}

任意の $N$ について、層共鳴のディリクレ核解析がすべての格子点で可去特異性条件を満たし、見かけの発散を防ぐ。したがって上限 $|J| \leq 1$ はすべての $N$ で一様に成立する。$N$ が増大すると、到達可能な $J$ 値の離散集合は区間 $[-1,1]$ 内で稠密になり、各有限段階で厳密な大域上限を保ちつつ連続理論を回復する。

これらの明示的列挙は、抽象的な有界性定理の具体的検証を与え、中程度のコンパクト化パラメータに対するディオファントス証明の数値的検査の基礎となる。

\section[ベイカー--キャンベル--ハウスドルフ展開とスピン被覆]{ベイカー--キャンベル--ハウスドルフ展開\\とスピン被覆の詳細}

\subsection{ローター冪に対するベイカー--キャンベル--ハウスドルフ公式}

フェルマーの最終定理の証明およびねじれ不整合の増幅では、ローターを整数乗に繰り返し持ち上げる。$\Omega = \exp(X)$ とし、$X \in \mathfrak{so}(3,1) \oplus \mathfrak{so}(3,1)$ をねじれ不整合二重ベクトルとする。すると
\[
\Omega^p = \exp(pX + \tfrac{p(p-1)}{2}[X,X] + \tfrac{p(p-1)(2p-1)}{12}[X,[X,X]] + \cdots).
\]
二次まで打ち切る（小さい不整合、または離散格子により高次交換子が抑制される場合に有効）と
\[
\log(\Omega^p) = pX + \tfrac{p(p-1)}{2}[X,X] + O(p^3).
\]
$J$ を定義する二次形式へ代入すると、主項の増幅
\[
J(\Omega^p) = p^2 J(\Omega) + O(p^3 \cdot \text{高次の入れ子交換子})
\]
が得られる。$p \geq 3$ では係数 $p^2$ が有限トーラス上のいかなる高次補正も厳密に支配し、第5章で用いた矛盾論証を正当化する。

\subsection{スピン被覆と主枝の制限}

群 $\mathrm{Spin}^+(3,1) \oplus \mathrm{Spin}^+(3,1)$ は固有時計ローレンツ群の二重被覆である。したがって指数写像
\[
\exp : \mathfrak{so}(3,1) \oplus \mathfrak{so}(3,1) \to \mathrm{Spin}^+(3,1) \oplus \mathrm{Spin}^+(3,1)
\]
は局所微分同相だが大域的には一対一ではない。核は元 $\pm 1$ からなる。したがって対数の主枝に注目し、
\[
U = \{ g \in \mathrm{Spin}^+(3,1) \oplus \mathrm{Spin}^+(3,1) \mid \| \log g \| < 2\pi \},
\]
と定義する。ここでノルムはキリング形式により誘導される。この近傍内で対数は一価かつ解析的である。

双対写像 $X \mapsto Xi$ はねじれ不整合二重ベクトルの係数に
\[
\alpha_a + \beta_a = \delta_a, \qquad |\delta_a| \leq \frac{\pi}{2}
\]
を主枝内で強制する。この上限は、双対セクターにおける完全 $2\pi$ 回転がスピン群の元 $-1$ に対応し、主領域の境界上に位置することから従う。これより大きな偏差は枝切りを横断することを要し、隣接粒子が連続ヌル並進子のみで差をとる大域的整合性条件により排除される。

\subsection{離散トーラス上における打ち切りの正当化}

有限トーラス $(\mathbb{Z}/N\mathbb{Z})^6$ 上ではすべてのラピディティが $2\pi$ 以下に有界である。したがってねじれ不整合二重ベクトル $X$ は $\|X\| \leq 3\pi$ を満たす。任意の固定コンパクト化パラメータ $N$ について、ベイカー--キャンベル--ハウスドルフ級数の高次項は一様に有界である。二次打ち切りによる誤差はしたがって $N$ のみに依存する定数で制御され、連続極限で消える。これがディオファントス証明で用いたすべての近似を正当化する。

ベイカー--キャンベル--ハウスドルフ展開とスピン被覆制限の組み合わせは、増幅論証と $J$ の有界性に対する完全に厳密な解析的基礎を与える。

\section{宇宙際タイヒミュラー理論対応表}

離散双四元数枠組みは、宇宙際タイヒミュラー理論の主要対象との精密な辞書を認める。次の表が主要な対応を要約する。

\begin{center}
\begingroup
\renewcommand{\arraystretch}{1.5}
\begin{tabular}{@{}p{6.5cm}p{6.5cm}@{}}
\toprule
\textbf{二重時空理論} & \textbf{宇宙際タイヒミュラー理論} \\
\midrule
ダガー作用 $R_{\rm usual}^\dagger$ & ログ・シータ解体（ログリンク） \\
双対ローター $R_{\rm dual}$ による乗法 & 再構成ステップ（$\Theta$-リンク） \\
ねじれ不整合スカラー $J$ & ログ体積／高さ関数 \\
離散双四元数環の有限単位群 & エタール・シータ関数／基体の有限単位 \\
離散トーラス $(\mathbb{Z}/N\mathbb{Z})^6$ & ログ・シータ格子の基本領域 \\
ヌル並進子 $T(a) = \frac12 a^\mu N_\mu$ & 算術データ／フロベニオイド射 \\
ローター冪の BCH 展開 & シータ関数の対数微分 \\
$\log \Omega$ の主枝 & $p$ 進対数の枝の選択 \\
ねじれ有界性 $|J| \leq 1$ & 基本領域上のログ体積の有界性 \\
IUT 型下降手続き & 宇宙際タイヒミュラー下降 \\
\bottomrule
\end{tabular}
\endgroup
\end{center}

これらの対応により、第5--8章のディオファントス証明の各ステップを宇宙際タイヒミュラー理論の言語へ直接翻訳でき、結果の独立検証と両枠組み間のさらなる相互発展への経路を開く。

\section{補助補題の証明}

\subsection{有限単位群に関する補題}

\begin{lemma}
コンパクト化パラメータ $N$ の離散二重時空に関連する整数双四元数環の単位群は有限である。
\end{lemma}

\begin{proof}
離散トーラス $(\mathbb{Z}/N\mathbb{Z})^6$ は有限集合である。リー代数からスピン群への指数写像は連続であり、連続写像の有限集合の像は有限である。したがって離散ラピディティパラメータから生じる異なるローターは有限個のみである。単位群はちょうどこれらのローターの集合（整数係数の作用を除いて）であるから、有限である。
\end{proof}

\subsection{ディリクレ核の可去特異性に関する補題}

\begin{lemma}
任意の整数 $N \geq 2$ および任意の整数 $k$ について、ディリクレ核は
\[
\lim_{\delta\theta \to 2\pi k} D_N(\delta\theta) = \pm 1.
\]
を満たす。
\end{lemma}

\begin{proof}
ディリクレ核は
\[
D_N(\theta) = \frac{\sin(N\theta/2)}{N\sin(\theta/2)}
\]
で与えられる。$\theta = 2\pi k$ では分子・分母ともにゼロとなり、$0/0$ の不定形を生じる。ロピタルの定理を一度適用すると
\[
\lim_{\theta \to 2\pi k} D_N(\theta) = \frac{(N/2)\cos(N\theta/2)}{(N/2)\cos(\theta/2)} \Big|_{\theta=2\pi k} = \cos(N\pi k) = (\pm 1)^N = \pm 1,
\]
となり、$N$ に依らない。したがって見かけの特異点はすべて可去である。
\end{proof}

\subsection{ローター冪下のねじれ増幅に関する補題}

\begin{lemma}
$\Omega = \exp(X)$、$X \in \mathfrak{so}(3,1) \oplus \mathfrak{so}(3,1)$ とする。すると
\[
J(\Omega^p) = p^2 J(\Omega) + O(p^3),
\]
ここで暗黙の定数はコンパクト化パラメータ $N$ のみに依存する。
\end{lemma}

\begin{proof}
ベイカー--キャンベル--ハウスドルフ公式を二次まで適用する：
\[
\log(\Omega^p) = pX + \tfrac{p(p-1)}{2}[X,X] + O(p^3).
\]
$J$ を定義する二次形式は次数2で斉次である。したがって主項はちょうど $p^2 J(X)$ を寄与する。すべての高次交換子は有限トーラス上で $N$ のみに依存する定数で有界される。ゆえに誤差項は $O(p^3)$ である。
\end{proof}

\subsection{主枝上限に関する補題}

\begin{lemma}
$\mathrm{Spin}^+(3,1) \oplus \mathrm{Spin}^+(3,1)$ 上の対数の主枝内部では、不整合係数は $|\delta_a| \leq \pi/2$ を満たす。
\end{lemma}

\begin{proof}
双対写像 $X \mapsto Xi$ は、双対セクターにおける完全 $2\pi$ 回転をスピン群の中心元 $-1$ と同一視する。主枝は対数のノルムが厳密に $2\pi$ 未満となる開集合として定義される。偏差 $|\delta_a| > \pi/2$ は総不整合をこの近傍外へ押し出し、主枝の選択と矛盾する。したがって $|\delta_a| \leq \pi/2$ である。
\end{proof}

これらの補題が第3--8章のすべての論証の技術的基礎を与える。

\end{document}
